% Created 2026-09-15 Tue 11:02
% Intended LaTeX compiler: pdflatex
\documentclass[11pt]{article}
\usepackage[utf8]{inputenc}
\usepackage[T1]{fontenc}
\usepackage{graphicx}
\usepackage{longtable}
\usepackage{wrapfig}
\usepackage{rotating}
\usepackage[normalem]{ulem}
\usepackage{amsmath}
\usepackage{amssymb}
\usepackage{capt-of}
\usepackage{hyperref}
\usepackage{listings}
\author{conn}
\date{\today}
\title{}
\hypersetup{
 pdfauthor={conn},
 pdftitle={},
 pdfkeywords={},
 pdfsubject={},
 pdfcreator={Emacs 29.1 (Org mode 9.6.6)}, 
 pdflang={English}}
\begin{document}

\tableofcontents

\section{Problem statement}
\label{sec:org8a92bff}
The goal is a Chapel program that computes the Euclidean minimum M(K) of
a number field K, for arbitrary degree, using Pierre Lezowski's
algorithm (src//.c) as the mathematical reference, but doing the numeric
heavy-lifting in floating point/Chapel instead of Pari, with Pari used
only for one-time field setup (already done by \texttt{generate\_field.gp}). The
current Chapel attempts (\texttt{firstCut.chpl} \ldots{} \texttt{sixthCut.chpl},
\texttt{kernel.chpl}, \texttt{kernel.py}, \texttt{exactMin.chpl}) frequently return wrong
minima and are hard-capped at degree ≤ 4. \#\# Diagnosis of the current
code Reading \texttt{chplPlan.org} and \texttt{note.txt} alongside the actual Chapel
sources shows the previous implementation attempts did \textbf{not} port
Lezowski's algorithm --- they invented an unrelated “AMR box/disk sieve”
that is unsound: / \texttt{kernel.chpl:54} (\texttt{testBoxCovering}) only checks
\texttt{N(x) < bound} and \texttt{N(unit·x) < bound}. Since every unit has
\texttt{|N(unit)| = 1}, the “unit” check is numerically almost the same test as
the first one --- it does nothing useful. Real absorption requires
testing \texttt{N(x - γ) < bound} against many candidate integers γ, which
never happens here (only γ = 0 is ever tried). * \texttt{sixthCut.chpl:68}
patches this by brute-force scanning integer offsets \texttt{\{-1,0,1\}} per
coordinate and unit exponents \texttt{\{-2..2\}} (hard-coded up to 3 coefficients
and 2 units --- this is the literal source of the “degree ≤ 4” ceiling).
This is an ad hoc heuristic, not derived from the field's geometry, so
it silently misses the real absorbing elements for larger
fields/discriminants/rank ≥ 2 unit groups --- exactly the failure the
user already reproduced for \texttt{x\textasciicircum{}3-3x-1} in \texttt{note.txt}. *
\texttt{exactMin.chpl:38,45} calls \texttt{kernel.testBoxCovering(currentBox, ...)}
where \texttt{currentBox} is a \texttt{Box} record, but \texttt{kernel.chpl}'s
\texttt{testBoxCovering} expects a plain \texttt{[] real(64)} --- these files are
already mutually inconsistent. * There is no equivalent anywhere of
\texttt{small\_elts} (\texttt{src/main.c:130}, enumerates candidate algebraic integers
bounding a given norm target) or of the real absorption formula
\texttt{norme\_bricolee} (\texttt{src/basef.c:439}). These are the mathematical core of
Lezowski's method, and both are missing entirely. * No unit-recentering
(\texttt{vecteurs\_possibles2=/=test\_des\_unites}, \texttt{src/main.c:676,757}), no
problem graph/cycle detection (Tarjan port, \texttt{src/graph.c:58,102}), and
no exact finalization via Pari (\texttt{calcul\_min\_liste\_cycles},
\texttt{src/pari\_min\_c.c:675}) exist in Chapel --- so even where the sieve
reports “uncleared”, there is no way to certify a matching lower bound,
which is why runs end in “uncertain” territory (\texttt{note.txt}).
\textbf{Conclusion: this needs a rewrite of the sieve/kernel, not a patch.} The
good news is the parts that are salvageable are salvageable: the
field\textsubscript{data.txt} format, \texttt{generate\_field.gp}, and the box→embedding
projection math already in \texttt{sixthCut.chpl}
(\texttt{computeBoxProjections=/=calculateBoxMaxNorm}, lines 30-65) are
structurally correct primitives --- they were just never wired to a
correct absorption test. \#\# Key algorithmic facts established during
research * Degree n, signature (r1,r2) always satisfy n = r1+2·r2; boxes
are naturally described in \emph{integer-coefficient space} (dimension n),
then projected into the r1+r2 real/complex embedding slots for the norm
bound --- this is exactly what the current complex-compact
\texttt{field\_data.txt} format and \texttt{computeBoxProjections} already do, so \textbf{no
changes to \texttt{generate\_field.gp=/=field\_data.txt} are needed for the
corrected core sieve}. * Lezowski's absorption test
(\texttt{test\_absorption=/=norme\_bricolee}) is sound using \textbf{only} the
enumerated small-integer candidate list per target bound K --- the
fundamental-unit action (\texttt{test\_des\_unites}) is an \emph{acceleration and
exact-certification} device, not a soundness requirement for the basic
covering test. This means a first, fully correct implementation can skip
unit-recentering and the cycle/graph machinery entirely and still be
mathematically rigorous for proving upper bounds. * A sieve that proves
“all boxes absorbed for bound K” only ever certifies \texttt{M(K) < K} (a
numeric upper bound). It can never certify a matching lower bound by
itself --- Lezowski's C code gets exact answers via the unit-orbit cycle
decomposition + an exact Pari evaluation at the cycle's rational
critical point. We can get the same certification much more simply: take
the box(es) that most resist absorption at maximum sieve depth (rational
coordinates), and make \textbf{one exact Pari call} to compute \texttt{min\_\{γ\} N(x-γ)}
there exactly. If that exact value matches the numeric upper bound from
the sieve, the minimum is proved exactly, without porting
Tarjan/graph-simplification at all. * Chapel domains must have a
compile-time-constant rank, so the arbitrary-degree candidate-integer
enumeration (\texttt{small\_elts}) cannot use a rank-n domain; it must use a
flattened 1-D index space over \texttt{(limsup-liminf+1)\textasciicircum{}degree} decoded via
mixed-radix arithmetic (generalizing the C code's manual odometer-style
increment in \texttt{main.c:150-232}), so it parallelizes cleanly with a single
\texttt{forall}. * Environment check: \texttt{gp} (Pari/GP), \texttt{gcc}, and \texttt{chpl} (2.9.0)
are all installed locally, and the existing \texttt{Makefile=/=src/*.c} can
build the reference \texttt{euclid} C binary --- this gives a ground-truth
oracle for validation without re-deriving expected minima by hand. \#
Proposed approach: staged rewrite \#\# Phase 1 --- Correct,
arbitrary-degree decision sieve (numeric upper bounds) New Chapel
modules (replacing the box-sieve logic in
\texttt{sixthCut.chpl=/=exactMin.chpl}; the field-file parser in \texttt{kernel.chpl}
is reused/cleaned up): * \texttt{NumberField.chpl}: parses \texttt{field\_data.txt}
into a record (degree, r1, r2, basis embeddings, unit embeddings) ---
cleanup of \texttt{kernel.chpl}'s loader. * \texttt{SmallElements.chpl}: implements
the arbitrary-degree analogue of \texttt{small\_elts} (\texttt{main.c:130}) ---
computes the per-coordinate bound from row-sums of the embedding matrix
and \texttt{K\textasciicircum{}\{1/degree\}}, then enumerates all integer coefficient vectors in
that box via the flattened mixed-radix \texttt{forall} described above, keeping
only those whose embedding actually falls in range (mirrors
\texttt{main.c:155-232}). * \texttt{Sieve.chpl}: keeps the \texttt{Box} record and
\texttt{computeBoxProjections=/=calculateBoxMaxNorm} primitives (ported from
\texttt{sixthCut.chpl:10-65}), and adds the real absorption test: for a box,
iterate over the current \texttt{SmallElements} list, subtract each candidate's
embedding from the box's projected center, and check
\texttt{calculateBoxMaxNorm(shifted) < K} (the correct generalization of
\texttt{norme\_bricolee}, \texttt{basef.c:439}) --- short-circuiting on the first
success and moving that candidate to the front of the list for locality
(mirrors \texttt{ramene\_en\_tete3}, \texttt{basef.c:500}). Boxes failing absorption are
bisected into \texttt{2\textasciicircum{}degree} children (this part of the existing AMR code,
\texttt{sixthCut.chpl:241-260}, is structurally fine and is kept), tracked with
a \texttt{list=/array per depth level and processed with =forall}. * A driver
that runs cut+absorb to a resolution limit for a given K, reports
success/failure, and performs bisection search over K, replacing the
naive \texttt{epsilon}-only binary search in \texttt{sixthCut.chpl:273-301} with the
same “shrink K based on evidence, keep best success” strategy as
\texttt{main.c}'s \texttt{main\_loop=/outer =while} (\texttt{main.c:2682} onward) generalized
to arbitrary degree. This phase alone fixes both reported defects: it is
mathematically sound (real absorption against real candidates) and has
no dimension-specific hard-coding, so it works for arbitrary degree. \#\#
Phase 2 --- Unit-action acceleration (optional but recommended before
scaling up) * Extend \texttt{NumberField.chpl} to compute the inverse embedding
matrix (\texttt{sigma\_inv}) via Chapel's \texttt{LinearAlgebra} module (unpack the
complex-compact embedding matrix to a real n×n matrix, invert, done once
at load time --- no \texttt{generate\_field.gp} changes needed). * Port
\texttt{vecteurs\_possibles2=/=test\_des\_unites} (\texttt{main.c:676-897}): apply each
fundamental unit to a surviving box's embeddings, use \texttt{sigma\_inv} to
re-round the twisted box back toward the origin in coefficient space,
and re-test absorption. This mirrors the real acceleration loop
(\texttt{boucle\_unites\_et\_decoupage\_v4}, \texttt{main.c:951}) and should sharply cut
down on the bisection depth needed for fields with larger regulators,
without which some fields may be numerically slow (not wrong, just slow)
under Phase 1 alone. \#\# Phase 3 --- Exact certification (targeted,
simpler than the C tool's graph/cycle machinery) * \texttt{Certify.chpl}: once
the Phase 1/2 sieve stalls at maximum depth for a shrinking K, collect
the most-resistant box center(s) (rational coordinates from the grid),
and shell out \textbf{once} to \texttt{gp} (Chapel \texttt{Subprocess} module) with a small
script that computes \texttt{min\_\{γ small\} N(x-γ)} exactly at that rational
point using Pari's exact arithmetic (reusing the same bounding-box logic
as \texttt{SmallElements}, just done exactly instead of in float). * Compare
the exact value to the numeric upper bound from the sieve: if they match
(typically true for the well-known small examples --- e.g. \texttt{x\textasciicircum{}3-3x-1} →
1/3), report the proven exact Euclidean minimum. This sidesteps porting
Tarjan (\texttt{graph.c:58}) and the problem-graph simplification
(\texttt{main.c:1345-1843}) entirely. * If this simplified certification is
ever insufficient (fields with genuinely non-isolated resistant orbits
under rank ≥ 2 unit groups), the full cycle-decomposition port remains a
documented fallback but is out of scope unless Phase 3's simpler
approach proves inadequate in testing. \#\# Parallelization roadmap (kept
in mind from Phase 1 onward, not built until correctness is solid) *
Single-locale multi-core: \texttt{forall} over the small-elements enumeration
and over each level's box list (already the pattern in the existing code
--- kept). * Multi-locale: shard the per-level “problems” list across
locales (e.g. \texttt{Block}-distributed array or a per-locale list merged with
an all-to-all at the end of each depth level), since each box's
absorption test is independent given the (replicated, typically small)
candidate list. * GPU:
\texttt{computeBoxProjections=/=calculateBoxMaxNorm=/absorption-against-candidate-list
is pure flat numeric work with no dynamic allocation in the inner loop
-{}-{}- a natural =foreach=/=@assertOnGpu} target once box/candidate data is
restructured into flat arrays-of-structs-of-arrays instead of
\texttt{list(Box)}. \# Validation strategy * Build the reference \texttt{src} C
\texttt{euclid} tool (\texttt{make} in \texttt{euclid-1.0/}, pari + gcc already present) to
generate ground-truth minima for a battery of fields spanning
degree/signature variety: the existing saved fixtures \texttt{q2.txt} (x²−2),
\texttt{cc23.txt} (x³+x²−1, disc −23), plus \texttt{x\textasciicircum{}3-3x-1} (rank-2 unit group,
known min 1/3, the case that broke the old sieve), \texttt{x\textasciicircum{}2-61}, and at
least one quintic/sextic example to prove the degree generalization. *
Add a small test harness script that runs \texttt{generate\_field.gp} (already
present) then the new Chapel pipeline for each fixture and diffs against
the C tool's reported minimum within tolerance. * Move \texttt{firstCut.chpl}
\ldots{} \texttt{sixthCut.chpl}, \texttt{exactMin.chpl}, \texttt{kernel.py}, \texttt{tmp.chpl} into an
\texttt{attic/} subdirectory rather than deleting them, since they're
superseded but harmless as history.
\end{document}